\documentclass[11pt,a4paper]{article}

% --- packages -----------------------------------------------------------------
\usepackage[utf8]{inputenc}
\usepackage[T1]{fontenc}
\usepackage{lmodern}
\usepackage[a4paper,margin=2.4cm]{geometry}
\usepackage{microtype}
\usepackage{xcolor}
\usepackage{listings}
\usepackage{hyperref}
\usepackage{enumitem}
\usepackage{tabularx}
\usepackage{booktabs}
\usepackage{longtable}
\usepackage{titlesec}
\usepackage{fancyhdr}
\usepackage{parskip}
\usepackage{titling}

% --- colours ------------------------------------------------------------------
\definecolor{tnbg}{HTML}{1A1B26}
\definecolor{tnfg}{HTML}{C0CAF5}
\definecolor{tnblue}{HTML}{7AA2F7}
\definecolor{tnred}{HTML}{F7768E}
\definecolor{tnyellow}{HTML}{E0AF68}
\definecolor{tngreen}{HTML}{73DACA}
\definecolor{tncomment}{HTML}{565F89}
\definecolor{linkcolor}{HTML}{2C5AA0}
\definecolor{lightgray}{HTML}{EAEAEA}

% --- hyperref -----------------------------------------------------------------
\hypersetup{
  colorlinks=true,
  linkcolor=linkcolor,
  urlcolor=linkcolor,
  citecolor=linkcolor,
  pdftitle={sotX v0.2.0 -- a verified microkernel BSD personality with deception, SMF, Crossbow and FMA},
  pdfauthor={Lucas Sotomayor (BASE4 Security)},
  pdfsubject={Microkernel operating system architecture and proposal},
  pdfkeywords={microkernel, Rust, capability-based security, deception, SMF, Crossbow, FMA, Wayland, formal verification, TLA+}
}

% --- listings -----------------------------------------------------------------
\lstdefinelanguage{rustlang}{
  keywords={fn,let,mut,pub,struct,enum,impl,trait,if,else,for,while,loop,match,return,use,mod,as,where,self,Self,const,static,unsafe,extern,move,async,await,box,break,continue,crate,dyn,false,true,in,ref,type},
  keywordstyle=\color{tnblue}\bfseries,
  ndkeywords={u8,u16,u32,u64,usize,i8,i16,i32,i64,isize,bool,char,str,Option,Result,Vec,Box,String},
  ndkeywordstyle=\color{tngreen}\bfseries,
  identifierstyle=\color{black},
  sensitive=true,
  comment=[l]{//},
  morecomment=[s]{/*}{*/},
  commentstyle=\color{tncomment}\itshape,
  stringstyle=\color{tnred},
  morestring=[b]",
  morestring=[b]',
}
\lstset{
  language=rustlang,
  basicstyle=\ttfamily\footnotesize,
  numbers=left,
  numberstyle=\tiny\color{tncomment},
  numbersep=8pt,
  frame=single,
  framesep=4pt,
  rulecolor=\color{lightgray},
  backgroundcolor=\color{lightgray!30},
  showstringspaces=false,
  breaklines=true,
  tabsize=2,
  captionpos=b,
}

% --- titlesec -----------------------------------------------------------------
\titleformat{\section}
  {\normalfont\Large\bfseries\color{tnblue}}
  {\thesection}{1em}{}
\titleformat{\subsection}
  {\normalfont\large\bfseries\color{tnblue!80!black}}
  {\thesubsection}{1em}{}

% --- header / footer ----------------------------------------------------------
\pagestyle{fancy}
\fancyhf{}
\fancyhead[L]{\textit{sotX v0.2.0 whitepaper}}
\fancyhead[R]{\textit{2026-04-08}}
\fancyfoot[C]{\thepage}
\renewcommand{\headrulewidth}{0.3pt}

% --- title metadata -----------------------------------------------------------
\title{
  \vspace{-1.5cm}
  \textbf{sotX v0.2.0} \\[0.4em]
  \large A Verified Microkernel BSD Personality with \\
  Deception, SMF, Project Crossbow and FMA \\[0.6em]
  \normalsize Architecture, Implementation, and Proposal
}
\author{
  Lucas Sotomayor \\
  \texttt{lucas@base4.io} \\
  BASE4 Security
}
\date{2026-04-08 \\ Codename: Project STYX}

% =============================================================================
\begin{document}
% =============================================================================

\maketitle
\thispagestyle{empty}

\begin{abstract}
\noindent
\textbf{sotX} is a verified microkernel operating system written from
scratch in Rust, layered with a SOT (Secure Object) exokernel, a NetBSD
\texttt{rump} kernel personality, a Linux ABI personality, a deception
subsystem with capability interposition and live migration, and a v0.2.0+
``Illumos Heirlooms'' wave that adds a Solaris-style Service Management
Facility (SMF), Project Crossbow virtual networking, and a Fault
Management Architecture (FMA) predictive self-healing engine. The
codebase is 108k+ LOC of Rust across 69 crates, 4.4k LOC of necessary
C glue (Limine + LKL/rump shims), 6 TLA+ specifications under TLC model
checking, and 16 PASS markers verified end-to-end in a single QEMU boot.
This document describes the architecture, the threat model the system
addresses, the v0.2.0 deliverables, and the technical proposal for what
sotX enables that no other open-source operating system in 2026 does.
\end{abstract}

\tableofcontents
\newpage

% =============================================================================
\section{Executive Summary}
% =============================================================================

sotX is the BSD personality branch of \texttt{sotX}, a research
microkernel built from scratch over the past year. As of v0.2.0
(commit \texttt{eb137e3}, 2026-04-08) it ships:

\begin{itemize}[leftmargin=1.4em,itemsep=0.3em]
  \item A 5-primitive verified microkernel
        (scheduling, IPC, capabilities, IRQ virtualisation, frame
        allocation), $\sim$10\,k LOC, with capability tables, CDT,
        per-core ticket-locked run queues, work stealing, EDF
        deadline scheduling, multi-address-space processes, and SMP
        on 1--4+ cores.

  \item A SOT (Secure Object) exokernel layer with 12 syscalls
        (300--311) for transactional secure objects, single-tier and
        two-phase commit, capability interposition with proxy
        domains, and a per-CPU 4096-entry provenance ring.

  \item Real \textbf{NetBSD rump kernel} linked into a freestanding
        sotX service: \texttt{librump.a + librumpvfs.a +
        librumpfs\_ffs.a + librumpdev.a + librumpdev\_disk.a + ...}
        ($\sim$2 MiB) bound through 60 hand-rolled \texttt{rumpuser\_*}
        hypercalls in \texttt{rumpuser\_sot.c}.

  \item A \textbf{Linux ABI personality} (LUCAS) with full
        \texttt{fork/execve/waitpid}, glibc and musl static + dynamic
        binaries, vDSO trampolines, full \texttt{mmap}/\texttt{brk},
        signal delivery (sync + async via timer ISR), real
        \texttt{clone(CLONE\_THREAD)}, futex, TLS, and the
        \texttt{services/lucas-shell} command interpreter.

  \item A \textbf{deception subsystem} (\texttt{personality/common/
        deception}) with anomaly detection (\textit{credential theft
        after backdoor, reconnaissance sweep, data exfiltration,
        persistence install, staged payload, privilege escalation,
        anomalous fan-out}), Migration Orchestrator, and four
        deception profiles (Ubuntu webserver, CentOS database, IoT
        camera, Windows file server) wired through capability
        interposition.

  \item An \textbf{Ed25519 + SHA-256 boot chain}: every signed
        binary in the initrd is hashed and Ed25519-verified against
        a public key compiled into init via
        \texttt{services/init/src/sigkey\_generated.rs}.
        Pubkey-pinning rejects swap-with-self-signed attempts.

  \item The \textbf{v0.2.0+ Illumos Heirlooms wave}: SMF (Service
        Management Facility) with active respawn via a new
        \texttt{SYS\_THREAD\_NOTIFY (143)} syscall and a
        topological-cascade dependency graph; \textbf{Project
        Crossbow} virtual switch with O(1) VNIC provision/revoke,
        token-bucket rate limiting and locally-administered MAC
        generation; and \textbf{FMA} (Fault Management Architecture)
        with 6 prioritised correlation rules over disk-retry,
        ECC-correctable, ECC-uncorrectable, thermal and PCI-link
        events plus provenance-anomaly counters.

  \item A renewed \textbf{Wayland compositor} with Tokyo Night
        theme: layer-shell (\texttt{zwlr\_layer\_shell\_v1}),
        xdg-shell popups + positioner, fractional scaling stub,
        wl\_output v4, multi-shape cursor theme, anti-aliased
        \texttt{embedded-graphics} mono font, modern decorations
        with traffic-light buttons, vertical gradient wallpaper,
        per-rectangle damage tracking, animation framework with
        easing curves, and a feature-gated \texttt{tiny-skia}
        path-rendering backend.

  \item A \textbf{6-tier boot demo pipeline} that exercises real
        type-level code paths from the personality crates and
        gates on PASS markers. End-to-end verification reaches
        \textbf{16 PASS markers in a single boot} on commit
        \texttt{eb137e3}, with zero panics and zero triple faults.

  \item \textbf{Formal verification}: 6 TLA+ specifications
        (capabilities, IPC, scheduler, sot\_capabilities,
        sot\_provenance, sot\_transactions) with matching TLC
        model-checking \texttt{.cfg} files. \texttt{just tlc-mc}
        runs all six and reports per-spec state counts.

  \item \textbf{Production hardening}: a strict CI gate
        (\texttt{cargo fmt --check} + \texttt{cargo clippy -D
        warnings}) on the kernel and core libraries, nightly ABI
        fuzzing of every public syscall with random arguments,
        reproducible disk-image builds (\texttt{SOURCE\_DATE\_EPOCH}),
        a per-PR ABI-diff bot, automated GitHub Releases with
        Ed25519-signed SDK tarballs, and \texttt{cargo audit} +
        \texttt{cargo deny} on every push.
\end{itemize}

The proposition is straightforward: no open-source operating system
in 2026 combines a verified microkernel core, a real NetBSD rump
in-tree, capability-based deception with live migration, an Illumos-style
service / network / fault management triad, and a formally-checked SOT
exokernel layer. sotX is small enough (108\,k LOC of Rust) to be
auditable but rich enough to host real Linux binaries
(\texttt{busybox}, \texttt{git clone}, dynamic glibc) and to drive
Wayland clients with a modern theme.

% =============================================================================
\section{Threat Model and Motivation}
% =============================================================================

\subsection{What sotX defends against}

sotX targets two threat classes that conventional operating systems
handle poorly:

\subsubsection*{Class A: post-exploitation lateral movement and
persistence}

Once an attacker gains code execution inside an unprivileged process
on a Linux server, the standard playbook is well known: enumerate
\texttt{/etc/shadow}, escalate via SUID binaries, install a persistence
mechanism in \texttt{/etc/cron.d}, then pivot to neighbouring services
through the kernel network stack. Conventional defences (SELinux,
AppArmor, seccomp) are policy filters that the attacker either bypasses
once or works around with documented techniques.

sotX changes the contract. Every kernel object is accessed through a
capability, and \emph{capabilities can be interposed} after the fact.
The deception subsystem watches the per-CPU provenance ring, identifies
the canonical \texttt{CredentialTheftAfterBackdoor} pattern (write a
fake \texttt{/sbin/sshd}, then read \texttt{/etc/shadow}, then read
\texttt{/proc/version}), and live-migrates the offending domain into
a deception profile. The attacker is now talking to a fake
\texttt{/proc/version} that reports Ubuntu 22.04, a tarpit
\texttt{/etc/shadow}, and a fake OpenSSH banner --- with no observable
disconnect. Kill chains stall on the \emph{wrong target}.

Project Crossbow extends this: as soon as a domain is moved into
deception mode, the Crossbow virtual switch provisions a VNIC with a
spoofed MAC and a 10\,Mbps token-bucket rate limit. The attacker can
no longer exfiltrate or pivot to a peer at line rate; they are
syntactically isolated in a sandbox at O(1) provisioning cost.

\subsubsection*{Class B: silent storage / hardware degradation}

Conventional storage stacks (ext4, XFS, NTFS) treat disk read retries
as transient noise and surface them only after the underlying device
fails. ECC-correctable RAM events are visible to mcelog but not to the
storage layer. Thermal warnings and PCI-link degrades go to nobody.

The FMA module consumes the same provenance ring the deception
subsystem uses, plus a hardware-fault history of 64 entries, and
applies six prioritised correlation rules. \textbf{One uncorrectable
ECC} fires \texttt{MigrateNow}. \textbf{Five correctable ECC events at
the same address} fire \texttt{MigrateNow}. \textbf{Five disk-retry
events at the same LBA within a 1\,B-TSC window} fire
\texttt{MigrateNow}. \textbf{Thermal warning AND $>$10
provenance-anomaly events} fires \texttt{MigrateNow}. The decision is
fed into the Migration Orchestrator which evacuates Tier-1 / Tier-2
transactions off the failing region into a healthy frame
\emph{before} the silicon throws a fatal exception.

This is not magic. It is the same idea as Solaris FMA, applied to a
microkernel with capability semantics where ``move this transaction
to a different memory region'' is a finite, well-defined operation
instead of the open-ended hand-wave it tends to become on Linux.

\subsection{Why a microkernel for this}

The deception and FMA layers cannot trust the very kernel they are
trying to defend. On a monolithic kernel, an attacker who reaches
\texttt{ring~0} via a single bug owns the deception detector itself.
On sotX, the kernel implements \emph{five primitives} and 35-ish
syscalls, of which the deception machinery is entirely userspace
(\texttt{services/init/} and the personality crates). A bug in the
deception detector cannot escalate. A bug in the personality cannot
escalate. The TCB is the verified microkernel plus the boot chain
\emph{and nothing else}.

The cost of this contract is two: (a) every cross-domain interaction
goes through capability invocation rather than a direct function call,
and (b) the userspace VMM, network stack, file system, and personality
layers must be re-implemented or wrapped from upstream. We pay (a)
with capability cache and fast IPC, and we pay (b) with the
\texttt{rump-vfs} service (real NetBSD VFS, BSD filesystems and
device drivers) plus the 60-function \texttt{rumpuser\_sot.c}
hypercall layer.

% =============================================================================
\section{Architecture}
% =============================================================================

\subsection{The five-primitive microkernel}

\texttt{kernel/src/} hosts the entire Ring~0 surface. The five
primitives are:

\begin{enumerate}[leftmargin=1.6em,itemsep=0.2em]
  \item \textbf{Scheduling} (\texttt{kernel/src/sched/}). Per-core
        ticket-locked run queues with work stealing, 4-level priority
        MLQ (\texttt{realtime/high/normal/low}), EDF for real-time
        threads, per-thread CPU-tick and memory-page resource limits
        enforced in the scheduler and frame allocator. The new
        \texttt{SYS\_THREAD\_NOTIFY} hook fires from
        \texttt{exit\_current()} so SMF can detect dying services
        without polling.

  \item \textbf{IPC} (\texttt{kernel/src/ipc/}). Synchronous endpoints
        (L4-style), asynchronous channels (lock-free SPSC ring
        buffers), notifications (one-shot event signals), routing
        layer for cross-node delivery. Capability transfer on
        \texttt{send} is supported. \texttt{call\_timeout} (syscall
        135) lets a caller wake up after $N$ ticks if the callee
        never replies.

  \item \textbf{Capabilities} (\texttt{kernel/src/cap/}). The
        capability table is a fixed-size array indexed by 32-bit
        handles with generation counters (ABA-safe). The Capability
        Derivation Tree (CDT) records delegation parent/child links
        so revocation cascades atomically. Every kernel object
        (thread, AS, endpoint, channel, frame, SO) is accessed
        through a capability with explicit rights bits.

  \item \textbf{IRQ virtualisation}. The kernel exposes IRQs as
        capabilities. Userspace drivers register an \texttt{irq\_cap}
        and receive notifications via the standard endpoint
        mechanism. The kernel handles only the bare minimum (LAPIC
        timer, IPIs, fault delivery).

  \item \textbf{Frame allocation} (\texttt{kernel/src/mm/}). Bitmap
        physical-frame allocator (single-byte-per-frame metadata),
        slab allocator on top (for per-CPU caches), demand paging
        delegated to a userspace VMM via the global fault handler.
        Frame reference counting (\texttt{frame\_refcount\_inc}) is
        used by CoW fork.
\end{enumerate}

The whole kernel fits in $\sim$18\,k LOC of Rust including 1340
\texttt{unsafe} blocks (mostly in \texttt{arch/x86\_64/}, all carrying
\texttt{// SAFETY:} comments after PR \#66). The boot binary is around
1\,MiB.

\subsection{The SOT exokernel layer}

\texttt{kernel/src/sot/} adds 12 syscalls on top of the microkernel
core. They implement Secure Objects, transactional updates, and
capability interposition.

\begin{table}[h]
\centering
\small
\begin{tabularx}{\textwidth}{rlX}
\toprule
\# & Syscall & Purpose \\
\midrule
300 & \texttt{so\_create} & Create a Secure Object (file/dir/process/socket/credential type) \\
301 & \texttt{so\_invoke} & Method call on an SO; routed through interposition if applicable \\
302 & \texttt{so\_grant} & Grant SO access to a domain (with optional interposition policy) \\
303 & \texttt{so\_revoke} & Revoke an SO grant (with epoch barrier) \\
304 & \texttt{so\_observe} & Observe an SO without acquiring rights \\
305 & \texttt{sot\_domain\_create} & Create an isolation domain \\
306 & \texttt{sot\_domain\_enter} & Switch into a domain \\
307 & \texttt{sot\_channel\_create} & Create a typed SOT channel \\
308 & \texttt{tx\_begin} & Begin a transaction (\texttt{ReadOnly} / \texttt{SingleObject} / \texttt{MultiObject}) \\
309 & \texttt{tx\_commit} & Commit (single phase for tiers 0/1, second phase for tier 2) \\
310 & \texttt{tx\_abort} & Abort + rollback via WAL \\
311 & \texttt{tx\_prepare} & \textbf{Tier 5}: Tier 2 (MultiObject) prepare phase for two-phase commit \\
\midrule
260 & \texttt{provenance\_drain} & Pop entries from the per-CPU provenance ring into a userspace buffer \\
261 & \texttt{provenance\_emit} & Inject a synthetic provenance entry (used by attacker simulator + FMA test) \\
262 & \texttt{provenance\_stats} & Read ring metadata (length, drops, total pushed) \\
\midrule
143 & \textbf{\texttt{thread\_notify}} & \textbf{NEW (PR \#51)}: register a notify-cap to be signalled when a target tid exits \\
144 & \texttt{notify\_poll} & Non-blocking poll for a pending notification \\
\bottomrule
\end{tabularx}
\caption{The SOT and notify syscall surface, syscalls 143/144 added in the v0.2.0 Illumos Heirlooms wave for SMF active respawn.}
\end{table}

The kernel \texttt{ProvenanceEntry} is a \texttt{\#[repr(C)]} 48-byte
record containing \texttt{epoch}, \texttt{domain\_id},
\texttt{operation}, \texttt{so\_type}, \texttt{so\_id}, \texttt{version},
\texttt{tx\_id}, \texttt{timestamp}. Every \texttt{handle\_tx\_commit}
and \texttt{handle\_tx\_abort} emits a corresponding
\texttt{OP\_TX\_COMMIT} / \texttt{OP\_TX\_ABORT} entry with
\texttt{so\_type = SOTYPE\_TX\_EVENT}, so userspace observers can drive
ZFS / HAMMER2 snapshot managers from real kernel transaction events
without polling.

\subsection{KARL: per-boot ID randomisation}

\texttt{kernel/src/karl.rs} derives a 64-bit \texttt{boot\_seed} from
\texttt{RDTSC} at boot time and exposes per-pool offsets. The effect is
the same as OpenBSD's KARL kernel re-link for everything visible
through the SOT syscall ABI:

\begin{lstlisting}[caption={Per-boot KARL drift, observed across two cold boots.},captionpos=b]
Boot 1: boot_seed = 0x2895a574d42c28d7   tx_id starts at 3559653379
Boot 2: boot_seed = 0x726387273efe151f   tx_id starts at 1056833539
\end{lstlisting}

Attackers cannot count from a known baseline. KARL is wired into the
\texttt{Pool<T>} generation counters for tx, thread and capability
pools.

% =============================================================================
\section{The v0.2.0 Illumos Heirlooms Wave}
% =============================================================================

The v0.2.0 batch (codename: \textbf{Project STYX}) takes inspiration
from Illumos / OpenIndiana / Solaris and adds three major subsystems
plus a Wayland compositor renewal and a 4-axis hygiene pass. All 20
units shipped as PRs \#53--\#72 over 2026-04-07/08 and merged to
\texttt{sotX} HEAD \texttt{eb137e3}.

\subsection{SMF (Service Management Facility)}

\texttt{services/init/src/supervisor.rs} now hosts a 16-service
fixed-size dependency graph. Each service entry carries:

\begin{lstlisting}
pub struct SmfService {
    name: &'static [u8],
    deps: &'static [&'static [u8]],
    state: SmfState,           // Offline / Starting / Online / Failed / Maintenance
    restart_policy: RestartPolicy,  // Never / OnFailure / Always
    tid: Option<u64>,
    notify_cap: Option<u64>,
    respawn_count: u32,
}
\end{lstlisting}

When a registered service is spawned, \texttt{register\_smf()} allocates
a notification, calls \texttt{sys::thread\_notify(thread\_cap,
notify\_cap)} so the kernel signals the notification on thread exit,
and stores the notify-cap in the service slot. The
\texttt{poll\_and\_respawn()} loop polls each notification non-blockingly
and, on a death event, transitions the slot to
\texttt{SmfState::Failed}, walks all dependents in topological order,
and respawns each according to its restart policy. The
\textit{honeypot~DB $\to$ deception services} cycle from the proposal
is encoded as a concrete dependency edge in the test fixture.

\textbf{Boot verification}: a real \texttt{attacker} thread is spawned,
the kernel reports its death, the supervisor detects the death event,
respawns the attacker, and prints
\verb|=== SMF respawn: PASS ===|. The full cycle is observed
end-to-end in QEMU on every boot.

The new kernel hook in \texttt{kernel/src/sched/mod.rs::exit\_current()}
fires \texttt{on\_thread\_exit()} which walks a fixed-size 64-entry
\texttt{THREAD\_NOTIFY} table under \texttt{spin::Mutex} and signals
each matching notification. The double-check race pattern
(\texttt{register\_thread\_notify} re-checks liveness after writing
its slot, \texttt{on\_thread\_exit} races for the same flag) is
documented and proven correct in the PR. The timer-ISR resource-limit
kill path intentionally does \emph{not} call \texttt{on\_thread\_exit}
to avoid lock-order inversion (NOTIFICATIONS lock inside an interrupt
handler); services with strict resource quotas should use
\texttt{RestartPolicy::Always} so the next poll cycle catches them
via passive liveness check.

\subsection{Project Crossbow: virtual networking}

\texttt{libs/sot-crossbow/} introduces a fixed-size 32-port
\texttt{VirtualSwitch} with O(1) provision/revoke and a token-bucket
rate limiter:

\begin{lstlisting}
pub struct VirtualSwitch {
    ports: [Option<Vnic>; MAX_VSWITCH_PORTS],   // 32
    byte_counters: [u64; MAX_VSWITCH_PORTS],
    last_window_tsc: [u64; MAX_VSWITCH_PORTS],
}

impl VirtualSwitch {
    pub fn provision(&mut self, domain: DomainId, bw_kbps: u32) -> Result<u8, CrossbowError>;
    pub fn revoke(&mut self, port: u8);
    pub fn route(&mut self, src_port: u8, pkt_len: usize, now_tsc: u64) -> SwitchDecision;
    pub fn lookup_domain(&self, domain: DomainId) -> Option<u8>;
    pub fn vnic(&self, port: u8) -> Option<&Vnic>;
}
\end{lstlisting}

Each \texttt{Vnic} carries a locally-administered MAC derived from
\texttt{(domain\_id, tsc)} via a Knuth golden-ratio mix
(\texttt{0x9E37\_79B9\_7F4A\_7C15}), so adjacent domain IDs and
consecutive provisionings always produce visibly different MACs.
The first byte has the LAA bit set (\texttt{0x02}) and the multicast
bit clear.

The token bucket refills once per 1-second TSC window;
\texttt{bw\_kbps * 125 = bytes/sec}. For a 10 Mbps cap (the default
for sandboxed APTs) the bucket allows 1.25 MB/s.
\texttt{route()} returns \texttt{SwitchDecision::Forward(port)},
\texttt{Drop}, or \texttt{RateLimited}.

The crate is wired into \texttt{personality/bsd/network/src/pf.rs}.
When \texttt{PfInterposer::enable\_deception(domain)} flips a domain
into deception mode, it also calls
\texttt{crossbow.provision(domain, DEFAULT\_BW\_LIMIT\_KBPS)}.
\texttt{disable\_deception(domain)} revokes the matching port. The
result: an attacker that triggers the deception migration is
\emph{simultaneously} talking to a fake \texttt{/proc/version}
\emph{and} routed through a 10 Mbps spoofed-MAC VNIC. They can no
longer pivot at line rate.

\textbf{Test coverage}: 18 unit tests for provision exhaustion,
revoke + reuse, MAC uniqueness across all 32 ports, locally-administered
bit, rate-limit trip, refill window, and TSC override. Boot-time smoke
prints \verb|=== Crossbow: PASS ===|.

\subsection{FMA (Fault Management Architecture)}

\texttt{libs/sot-fma/} introduces a small predictive engine with
\texttt{FMA\_FAULT\_HISTORY = 64} hardware-fault entries and
\texttt{FMA\_PROVENANCE\_HISTORY = 256} anomaly-counter entries:

\begin{lstlisting}
#[derive(Clone, Copy, Debug, PartialEq, Eq)]
pub enum HwFaultClass {
    DiskReadRetry, EccCorrectable, EccUncorrectable,
    ThermalWarn, PciLinkDegrade,
}

#[derive(Clone, Copy, Debug, PartialEq, Eq)]
pub enum MigrationRecommendation {
    MigrateNow { reason: &'static str },
    MonitorClosely { reason: &'static str },
    NoAction,
}

pub struct FaultManagement { /* fixed-size, no_std, forbid(unsafe_code) */ }

impl FaultManagement {
    pub fn ingest_provenance(&mut self, entries: &[ProvenanceEntry]);
    pub fn ingest_hw_fault(&mut self, fault: HwFault);
    pub fn predict_failure(&self) -> MigrationRecommendation;
    pub fn fault_count(&self, class: HwFaultClass, window_tsc: u64, now_tsc: u64) -> usize;
}
\end{lstlisting}

The predictor (\texttt{libs/sot-fma/src/predictor.rs}) applies six
prioritised rules:

\begin{enumerate}[leftmargin=1.6em,itemsep=0.2em]
  \item Any single \texttt{EccUncorrectable} $\to$
        \texttt{MigrateNow \{reason: "ecc uncorrectable"\}}
  \item 5+ \texttt{EccCorrectable} at the same address $\to$
        \texttt{MigrateNow \{reason: "ecc soft errors clustering"\}}
  \item 5+ \texttt{DiskReadRetry} at the same location within
        a 1-billion-TSC window $\to$
        \texttt{MigrateNow \{reason: "disk read retries"\}}
  \item \texttt{ThermalWarn} AND \texttt{prov\_window\_anomalies > 10} $\to$
        \texttt{MigrateNow \{reason: "thermal + behavioral anomaly"\}}
  \item Any single fault $\to$ \texttt{MonitorClosely \{reason: ...\}}
  \item Otherwise $\to$ \texttt{NoAction}
\end{enumerate}

The 26 unit tests cover every threshold boundary, priority ordering,
ring rotation, and TSC saturation. Boot-time smoke
(\texttt{services/init/src/fma\_demo.rs}) synthesises fake disk
retries and uncorrectable ECC events and asserts the right
recommendation, printing \verb|=== FMA: PASS ===|.

\textbf{Known follow-up}: the local \texttt{ProvenanceEntry} mirror in
\texttt{libs/sot-fma/src/lib.rs} is 40 bytes (rustc default field
layout, no \texttt{\#[repr(C)]}) while the kernel's struct is 48 bytes
\texttt{\#[repr(C)]}. Fine for the synthesised demo today; needs to be
made layout-compatible before real drained-ring data is plumbed in.

% =============================================================================
\section{The Wayland Compositor Renewal}
% =============================================================================

\texttt{services/compositor/} went from a barebones engineering preview
(8x8 bitmap font, 12x12 white arrow cursor, hardcoded BGRA
\texttt{0xFF2D2D3D} background, no animations, no popups) to a
modern Wayland compositor with the Tokyo Night palette in twelve
parallel units (G1--G12).

\subsection{Visual upgrade summary}

\begin{table}[h]
\small
\centering
\begin{tabularx}{\textwidth}{lXl}
\toprule
\textbf{Was} & \textbf{Is (after v0.2.0)} & \textbf{PR} \\
\midrule
Solid \#2D2D3D background & Tokyo Night vertical gradient (\#1A1B26 $\to$ \#10101C), 8.6\,KiB row cache, optional BMP image tile path & \#57 \\
8$\times$8 bitmap font, 95 ASCII glyphs & \texttt{embedded-graphics FONT\_10X20} (large) + \texttt{FONT\_6X10} (compact), anti-aliased, with a \texttt{FbTarget} adapter implementing \texttt{DrawTarget} & \#56 \\
12$\times$12 white arrow cursor only & 8 cursor shapes (Default, Pointer, Text, Wait, Move, ResizeNS, ResizeEW, NotAllowed) with hand-crafted BGRA glyphs and hot-spots, in static \texttt{rodata} & \#55 \\
24\,px solid title bar, red square close button & 28\,px Tokyo Night title bar with three traffic-light circles (red \#F7768E, yellow \#E0AF68, green \#73DACA, 14\,px diameter), rounded corners & \#62 \\
Single \texttt{bool} damage flag, full-screen redraw & 32-rectangle dirty list with coalesce-on-overflow and fullscreen-escalation if union $>$ 3/4 of screen & \#53 \\
No popups (\texttt{xdg\_wm\_base} v2 advertised but only \texttt{xdg\_toplevel} implemented) & Full \texttt{xdg\_positioner} + \texttt{xdg\_popup} with all anchor / gravity combinations and on-the-fly geometry recompute on \texttt{reposition} & \#70 \\
No layer-shell, no status bars & \texttt{zwlr\_layer\_shell\_v1} v4 with 4-layer z-order (Background $\to$ Bottom $\to$ Top $\to$ Overlay), anchor bitmask, exclusive\_zone reservation, margin support, 8 layer-surface pool & \#69 \\
No \texttt{wl\_output} advertised & Single-output \texttt{wl\_output} v4 global with framebuffer geometry/mode/scale/name/description from BootInfo & \#61 \\
No HiDPI awareness & \texttt{wp\_fractional\_scale\_v1} stub, advertises 120/120 (1.0x) per surface, foundation for future fractional scaling & \#60 \\
No animations & 32-animation pool with \texttt{Easing::\{Linear, EaseOut, EaseInOut\}}, \texttt{interpolate\_color}, no \texttt{powf} (manual cubic) & \#58 \\
No vector rendering & \texttt{tiny-skia} path-rendering backend gated behind \texttt{skia} Cargo feature: \texttt{SkiaCanvas::fill\_rounded\_rect}, \texttt{stroke\_rounded\_rect}, \texttt{fill\_circle}, \texttt{linear\_gradient}, \texttt{present()} & \#54 \\
Single hardcoded compositor & New \texttt{services/sot-statusbar/} layer-shell client with full Wayland handshake, \texttt{ack\_configure} round-trip, Tokyo Night palette, clock + memory + procs sampling & \#72 \\
\bottomrule
\end{tabularx}
\caption{The Wayland compositor upgrade across the v0.2.0 G1--G12 wave.}
\end{table}

\subsection{Compositor architecture invariants}

The whole compositor is \texttt{\#![no\_std]} and \texttt{\#![no\_main]},
single-threaded, with all per-client state in fixed-size arrays:

\begin{itemize}[leftmargin=1.4em,itemsep=0.2em]
  \item \texttt{[WlClient; 8]} --- max 8 concurrent clients
  \item \texttt{[Surface; 32]} --- max 32 surfaces
  \item \texttt{[Toplevel; 16]} --- max 16 toplevel windows
  \item \texttt{[ShmPool; 16]}, \texttt{[ShmBuffer; 32]}
  \item \texttt{[LayerSurface; 8]} for layer-shell
  \item \texttt{[Positioner; 16]}, \texttt{[Popup; 16]} for xdg-popup
  \item \texttt{[FractionalScale; 32]} for fractional scaling
\end{itemize}

The compose loop reads damage rects, clears those regions with the
wallpaper, layouts layer surfaces, then draws in strict z-order:
Background $\to$ Bottom $\to$ toplevels $\to$ Top $\to$ cursor $\to$
Overlay. Toplevel render area is reduced by the cumulative
\texttt{exclusive\_zone} of all visible layer surfaces so toplevels
don't overlap a status bar.

% =============================================================================
\section{Boot Verification: 16 PASS Markers}
% =============================================================================

\texttt{just build \&\& just run} on commit \texttt{eb137e3} reaches
\textbf{16 PASS markers in a single boot} with zero panics:

\begin{lstlisting}[language={},caption={Live serial output (excerpt) from the post-v0.2.0 boot.}]
=== compositor input: WIRED ===                  (line 218)
signify: verified=22 failed=0 missing=0          (line 237)
=== Signify boot chain: PASS ===                 (line 238)
=== STYX TEST: SOT exokernel syscalls 300-310 ===
STYX TEST: ALL PASS                              (line 279)
=== posix-test: PASS ===                         (line 446)
=== kernel-test: PASS ===                        (line 566)
RUMP-VFS-TEST: PASS                              (line 751)
=== Tier 3 demo: PASS ===                        (line 806)
=== FMA: PASS ===                                (line 816)
=== Tier 4 demo: PASS ===                        (line 1409)
=== Crossbow: PASS ===                           (line 1417)
=== Tier 5 demo: PASS ===                        (line 16948)
=== abi-fuzz: 10000/10000 survived ===           (line 18702)
=== SMF respawn: PASS ===                        (line 19094)
DLTEST: PASS - dynamic linking works!            (line 19106)
DLTEST: mul PASS                                 (line 19108)
\end{lstlisting}

Each marker exercises real type-level code paths from the personality
crates, not stubs. \texttt{abi-fuzz} alone hammers 10\,000 random
syscall arguments without panicking the kernel. \texttt{Tier 5}
exercises a real two-phase commit, 5\,000 random SOT syscalls with
garbage args, concurrent transaction stress, and TCG-calibrated
performance numbers next to published seL4 / NOVA / Fiasco / Linux
references.

% =============================================================================
\section{Production Hardening}
% =============================================================================

\subsection{Strict CI gate}

The new \texttt{lint-strict-core} job in
\texttt{.github/workflows/ci.yml} runs \texttt{cargo fmt --check} and
\texttt{cargo clippy --target x86\_64-unknown-none -- -D warnings} on
\texttt{sotos-kernel}, \texttt{sotos-common}, \texttt{sot-fma},
\texttt{sot-crossbow}, and \texttt{sotos-ld} \emph{without}
\texttt{continue-on-error}. Eleven inline clippy fixes (mostly
\texttt{.flatten()}, \texttt{div\_ceil}, \texttt{iter\_mut}) plus
crate-level targeted \texttt{\#![allow(...)]} for known dead code with
per-lint justification comments. \emph{No} blanket
\texttt{\#![allow(warnings)]}.

A drive-by fix folds \texttt{.got} and \texttt{.got.plt} into
\texttt{.data} in \texttt{kernel/linker.ld}, resolving a pre-existing
\texttt{.text} / \texttt{.got} section overlap that was blocking
\texttt{just build} on fresh checkouts.

\subsection{Audit + supply chain}

A new \texttt{audit.yml} workflow runs \texttt{cargo audit} (RustSec
advisory database) and \texttt{cargo deny check} on every push and PR.
\texttt{deny.toml} declares an explicit license allowlist
(MIT, Apache-2.0, BSD-2/3, ISC, MPL-2.0, Unicode, Zlib, CC0, 0BSD), no
GPL/AGPL, sources restricted to \texttt{crates.io}, and yanked crates
denied. The job hard-fails on any RustSec advisory or license
violation.

\subsection{Reproducible builds}

\texttt{repro.yml} runs a Linux + Windows matrix that builds
\texttt{target/sotx.img} \emph{twice on each host} via
\texttt{python scripts/mkimage.py}, then compares the two SHA-256s.
A final \texttt{compare} job downloads both host artifacts and
asserts they're equal. Determinism is enforced by pinning all FAT32
directory-entry timestamps to \texttt{SOURCE\_DATE\_EPOCH} (default
\texttt{1700000000}, 2023-11-14) via UTC \texttt{gmtime} in
\texttt{mkimage.py}.

\subsection{ABI fuzzing and ABI diff bot}

\texttt{services/abi-fuzz/} is a no\_std binary with a deterministic
xorshift64 PRNG seeded by a fixed prime. It picks a random syscall
from a curated whitelist (excluding the obviously-blocking and
destructive numbers), fills six argument registers with random
\texttt{u64}s, invokes via the raw \texttt{syscall6} helper, and loops
for 10\,000 iterations. The kernel must not panic, regardless of arg
garbage. The harness prints \verb|=== abi-fuzz: 10000/10000 survived ===|
on success. \texttt{fuzz.yml} runs this nightly under QEMU.

\texttt{abi-diff.yml} parses the \texttt{Syscall} enum from
\texttt{libs/sotos-common/src/lib.rs} on the PR base and head, computes
the symmetric diff (added / removed / renumbered variants), and posts
a sticky markdown comment on the PR. The script uses only Python
stdlib so it runs on Linux and Windows. Currently the parser correctly
identifies all 81 \texttt{Syscall} variants in the live tree.

\subsection{Reproducible release pipeline}

\texttt{release.yml} triggers on \texttt{git tag v*}, installs nightly
Rust + \texttt{just} + QEMU, runs \texttt{just sigmanifest \&\& just
image}, asserts the 16 PASS markers in a 4-minute QEMU boot, runs
\texttt{bash scripts/make-sdk.sh --version \$\{TAG\}}, signs the
tarball with Ed25519 (key sourced from
\texttt{secrets.SIGNIFY\_KEY\_BYTES} as a base64-encoded 32-byte raw
private key, decoded into a temp file with \texttt{trap} cleanup),
and uploads \texttt{.tar.gz}, \texttt{.sha256} and \texttt{.sig} via
\texttt{softprops/action-gh-release@v2}. A repo guard
(\texttt{if: github.repository == 'sotomayorlucas/sotX'}) prevents
fork-tag accidents.

% =============================================================================
\section{Formal Verification}
% =============================================================================

\texttt{formal/} contains six TLA+ specifications and matching
\texttt{.cfg} files for TLC model checking. \texttt{just tlc-mc} runs
all six and reports per-spec state counts:

\begin{lstlisting}[language={}]
==> TLC  ok  capabilities      (79137 distinct states, depth 10)
==> TLC  ok  ipc               ( 1498 distinct states, depth  5)
==> TLC  ok  scheduler         (  514 distinct states, depth 18)
==> TLC  ok  sot_capabilities  (  850 distinct states, depth  4)
==> TLC  ok  sot_provenance    ( 1462 distinct states, depth 13)
==> TLC  ok  sot_transactions  (  602 distinct states, depth  9)
tlc: 6 passed, 0 failed, 0 skipped
\end{lstlisting}

The first run of TLC found and fixed real bugs in five specifications:
unbounded \texttt{CHOOSE} for NULL sentinels (anchored to one-element
sets), an over-strict \texttt{Rollback} invariant
(\texttt{sot\_transactions::Rollback} required \emph{every} WAL entry's
\texttt{oldVal} to match \texttt{soValue}, but \texttt{RestoreFromWal}
picks the \emph{earliest} per object), stale WAL after abort,
duplicate-thread \texttt{Enqueue} (\texttt{scheduler::Enqueue} allowed
the same thread twice in a CPU's runqueue), and a looping
\texttt{RECURSIVE Descendants} in \texttt{capabilities::Revoke}
(replaced with an explicit fixed-point iteration bounded by
\texttt{MaxCaps + 1}).

One open finding is parked: \texttt{sot\_transactions::Atomicity} is
violated when Tier 1 \texttt{T1Write} and Tier 2 \texttt{T2Begin}
interleave on the same object; the spec's lock semantics don't honour
the \texttt{T2} prepared state across tiers. The \texttt{.cfg}
restricts checking to a single transaction so the structural
invariants still verify.

\texttt{just tlc} runs the SANY syntax/level checker as a faster
pre-check, and the \texttt{tla-tlc} CI job runs full model checking on
every push.

% =============================================================================
\section{Codebase Inventory}
% =============================================================================

As of v0.2.0 (\texttt{eb137e3}, 2026-04-08):

\begin{table}[h]
\centering
\small
\begin{tabular}{lrr}
\toprule
\textbf{Language} & \textbf{File count} & \textbf{LOC} \\
\midrule
Rust (.rs) & 376+ & $\sim$108\,k \\
C (.c, vendor + glue) & 13 & $\sim$4.4\,k \\
C++ & 0 & 0 \\
Python (.py) & 50+ & $\sim$5\,k \\
TLA+ (.tla) & 6 & $\sim$1.5\,k \\
Markdown (.md) & 16+ & $\sim$10\,k \\
TOML (.toml) & 105 & --- \\
Total & --- & $\sim$130\,k \\
\midrule
Crates in workspace & \multicolumn{2}{r}{69} \\
Kernel pools (static memory) & \multicolumn{2}{r}{$\sim$4.37\,MiB} \\
Largest single pool & \multicolumn{2}{r}{\texttt{FRAME\_REFCOUNT} 4\,MiB} \\
\texttt{unsafe} blocks (kernel) & \multicolumn{2}{r}{1340} \\
TLA+ specifications (TLC-checked) & \multicolumn{2}{r}{6} \\
PASS markers in a clean boot & \multicolumn{2}{r}{16} \\
\bottomrule
\end{tabular}
\caption{Codebase inventory after the v0.2.0 wave. C is justified
(Limine bootloader glue + LKL/rump shims for the in-tree NetBSD
\texttt{rump} layer); no migration needed.}
\end{table}

The C inventory is fully justified: \texttt{limine/limine.c} is the
external bootloader, \texttt{services/init/lkl/} and
\texttt{services/lkl-server/} are the LKL hypercall bridge for the
NetBSD \texttt{rump} layer (whose API is C and cannot be reasonably
re-implemented in Rust without rewriting the whole rump kernel). All
other userspace and kernel code is Rust. There is zero C++.

% =============================================================================
\section{What sotX enables}
% =============================================================================

The proposition for v0.2.0+ is short:

\begin{enumerate}[leftmargin=1.6em,itemsep=0.4em]
  \item \textbf{Run real Linux binaries on a verified microkernel.}
        Static and dynamic glibc + musl, \texttt{busybox sh}, real
        \texttt{git clone}, real \texttt{nano} editor, all under the
        LUCAS personality with the kernel TCB capped at 18\,k LOC.

  \item \textbf{Detect post-exploitation lateral movement and
        live-migrate the attacker into a deception profile} without
        breaking their TCP state. Anomaly $\to$ Migration Orchestrator
        $\to$ Crossbow VNIC sandbox at 10\,Mbps $\to$ fake
        \texttt{/proc/version}, fake \texttt{/sbin/sshd}, tarpit
        \texttt{/etc/shadow}.

  \item \textbf{Predict hardware failures from soft signals} (disk
        retries, ECC corrections, thermal warnings) and evacuate
        Tier-1 / Tier-2 transactions \emph{before} the silicon
        throws a fatal exception. The same provenance ring that
        feeds the deception subsystem feeds FMA.

  \item \textbf{Respawn dependent services in topological order} when
        a critical one dies, via SMF. The kernel signals
        \texttt{SYS\_THREAD\_NOTIFY}; the supervisor cascades through
        the dependency graph; the attacker's TCP connection to a
        deception service survives the restart.

  \item \textbf{Sign and verify the entire boot chain} with Ed25519,
        with pubkey-pinning to defeat swap-with-self-signed attacks.
        Tampering with one byte of any signed binary is caught at
        boot before the binary executes.

  \item \textbf{Reproduce the disk image bit-for-bit} across Linux
        and Windows hosts (\texttt{SOURCE\_DATE\_EPOCH}), so a release
        SHA-256 actually means something.

  \item \textbf{Drive a Wayland compositor with a modern Tokyo Night
        theme}, real anti-aliased fonts, multi-shape cursors,
        traffic-light decorations, layer-shell status bar, popups,
        animations, all in a \texttt{\#![no\_std]} userspace
        binary. The compositor talks to the kernel via the same
        capability-based IPC every other service uses.

  \item \textbf{Verify the kernel's safety properties} with TLA+
        model checking on six specifications, with the model
        checker running in CI on every push.
\end{enumerate}

% =============================================================================
\section{Roadmap and known follow-ups}
% =============================================================================

\subsection{Post-v0.2.0 follow-ups (non-blocking)}

\begin{itemize}[leftmargin=1.4em,itemsep=0.3em]
  \item \textbf{sot-statusbar happy path}: \texttt{init}'s
        \texttt{spawn\_process} passes \texttt{cap=0}, so the
        statusbar can't allocate caps for the SHM pool. Needs
        kernel-side spawn (like the compositor) or cap-forwarding
        via \texttt{bootinfo\_write}. The statusbar gracefully
        connects-and-waits today.

  \item \textbf{sot-fma \texttt{ProvenanceEntry} ABI compatibility}:
        local mirror is 40 bytes, kernel struct is 48 bytes
        \texttt{\#[repr(C)]}. Fix: import the upstream type or add
        \texttt{\#[repr(C)]} + a \texttt{const\_assert!} on size.

  \item \textbf{Multi-monitor real implementation}: \texttt{wl\_output}
        is currently a single-output stub. Hot-plug and multi-seat
        are deferred.

  \item \textbf{Real fractional scaling render path}: today the
        \texttt{wp\_fractional\_scale\_v1} stub advertises 1.0$\times$
        but the compositor renders at 1$\times$ regardless. HiDPI
        font rendering and surface scaling are deferred.

  \item \textbf{IME, clipboard, screen capture, gesture support,
        multi-seat}: real Wayland protocols, each a multi-day effort.

  \item \textbf{LTP (Linux Test Project) subset under LUCAS}:
        blocked on a host script to fetch and cross-build LTP.

  \item \textbf{WHPX/KVM IRQ model rework}: WHPX boots Limine but
        crashes on LAPIC injection
        (\texttt{Unexpected VP exit code 4}). The kernel's IRQ model
        isn't WHPX-compatible without a deeper rework. Today we use
        a TCG-calibration approach instead.

  \item \textbf{CHERI Morello target}: an
        \texttt{aarch64-morello-purecap} target plus an in-kernel
        CHERI enforcement path, so the software model in
        \texttt{services/sot-cheri} becomes a hardware-checked one.
\end{itemize}

\subsection{Planned for v0.3.0}

\begin{itemize}[leftmargin=1.4em,itemsep=0.3em]
  \item Real OpenZFS \texttt{libzpool} link (today the snapshot
        manager is in-memory). Connect to SOT block device caps and
        replace the demo plumbing.
  \item \texttt{rump\_init} curlwp rewrite to unblock the full
        \texttt{rump} kernel boot beyond the current $\sim$64
        hypercall mark.
  \item bhyve VT-x / WHPX support so guest VMs can run with
        introspectable caps.
  \item HAMMER2 real userspace library + clustering for
        distributed deception.
  \item GPU rendering: DRM/KMS direct, OpenGL ES, or Vulkan. Today
        \texttt{tiny-skia} gives software path-based rendering as
        the immediate win.
\end{itemize}

% =============================================================================
\section{Conclusion}
% =============================================================================

sotX v0.2.0 ships a verified microkernel BSD personality OS in
$\sim$108\,k LOC of Rust, with a SOT exokernel layer, a real NetBSD
\texttt{rump} kernel, a Linux ABI personality that runs real glibc
and musl binaries, capability-based deception with live migration,
SMF active-respawn supervision, Project Crossbow virtual networking,
FMA predictive self-healing, a Tokyo Night Wayland compositor, and a
production hardening pipeline that enforces formatting, lints,
licenses, advisories, ABI surface, reproducible builds and
boot-smoke testing on every PR. Six TLA+ specifications are model
checked in CI on every push.

This is, to the author's knowledge, the only operating system in 2026
that combines all of these in a single auditable codebase. The current
work was driven by an exploration of how Solaris, OpenIndiana, OpenBSD
and Illumos solved problems that mainstream Linux distributions still
hand-wave about, ported into a verified microkernel where the kernel
TCB is small enough to be meaningfully audited and the userspace is
modular enough that any single component can be replaced or
upgraded independently.

The next release (v0.3.0) will focus on closing the four big
follow-ups: real \texttt{libzpool}, \texttt{rump\_init} curlwp,
bhyve guest VMs, and HAMMER2 clustering for distributed deception.

\vspace{1em}
\hrule
\vspace{0.5em}
\noindent\textit{Repository:}
\href{https://github.com/sotomayorlucas/sotX}{\texttt{github.com/sotomayorlucas/sotX}}
\quad \textit{Branch:} \texttt{sotX}
\quad \textit{HEAD:} \texttt{eb137e3}
\quad \textit{Date:} 2026-04-08

\noindent\textit{Author:} Lucas Sotomayor, BASE4 Security
\quad \textit{Codename:} Project STYX

% =============================================================================
\end{document}
% =============================================================================
